\section{Chapitre 1 : Contexte Général et État de l'Art}
\addcontentsline{toc}{section}{Chapitre 1 : Contexte Général et État de l'Art}

L'histoire de l'éducation à distance a connu de multiples bouleversements paradigmatiques, évoluant des cours par correspondance acheminés par voie postale jusqu'aux écosystèmes d'apprentissage hautement interactifs que nous connaissons aujourd'hui. L'avènement des premiers réseaux informatiques dans les années quatre-vingt-dix a marqué la genèse des systèmes de gestion de l'apprentissage modernes, initialement conçus comme de simples dépôts de documents numérisés où les apprenants pouvaient télécharger des ressources textuelles. Rapidement, ces plateformes se sont enrichies de fonctionnalités de communication asynchrone, telles que les forums de discussion et les courriers électroniques internes, favorisant ainsi l'émergence des premières véritables communautés virtuelles d'apprentissage. Des solutions open-source emblématiques, à l'instar de Moodle, ou des plateformes propriétaires telles que Blackboard et Canvas, ont progressivement dominé le marché institutionnel en offrant une panoplie d'outils administratifs et pédagogiques. Ces outils ont permis aux enseignants de structurer leurs cours sous forme de modules chronologiques, de distribuer des devoirs, d'organiser des évaluations sous forme de questionnaires à choix multiples et de maintenir un registre centralisé des notes. Bien que ces systèmes traditionnels aient considérablement démocratisé l'accès au savoir, abolissant les frontières géographiques et temporelles, ils reposent fondamentalement sur un paradigme de transmission unidirectionnelle des connaissances, souvent qualifié d'approche "one-size-fits-all". Dans ce modèle canonique, chaque apprenant, indépendamment de son bagage cognitif initial, de sa vélocité d'assimilation ou de son style d'apprentissage préférentiel, est soumis exactement au même cheminement pédagogique. Cette rigidité structurelle constitue aujourd'hui l'un des freins majeurs à l'efficacité de l'enseignement en ligne, se traduisant empiriquement par des taux d'attrition préoccupants dans les cours ouverts massifs et un manque d'engagement intellectuel profond chez de nombreux étudiants.

Parallèlement à cette crise de la personnalisation pédagogique, la massification de l'évaluation à distance a mis en exergue une problématique de gouvernance majeure : la garantie de l'intégrité académique. Historiquement, l'évaluation des connaissances s'est toujours déroulée dans des amphithéâtres ou des salles de classe sous la surveillance directe et inquisitrice de personnels enseignants, garantissant ainsi que l'étudiant évalué est bien l'auteur des réponses fournies et qu'il ne bénéficie d'aucune assistance non autorisée. La transposition de ce processus d'évaluation dans des environnements domestiques non contrôlés a engendré une prolifération de comportements frauduleux d'une sophistication croissante. Les étudiants peuvent désormais recourir à l'usurpation d'identité en se faisant remplacer par des tiers, utiliser des dispositifs électroniques dissimulés pour consulter des bases de données externes, partager l'écran de leur ordinateur avec des collaborateurs à distance via des logiciels de bureau à distance, ou encore exploiter des failles techniques pour intercepter les questions d'examen. Face à ces menaces polymorphes qui sapent la valeur même des diplômes, les institutions éducatives se sont tournées vers des solutions de télésurveillance, ou proctoring. Les premières itérations de ces systèmes reposaient sur une surveillance humaine synchrone, où un surveillant distant observait plusieurs étudiants simultanément via le flux de leur webcam et le partage de leur écran. Bien que relativement efficace, cette méthode s'est avérée coûteuse, difficilement scalable pour de grandes cohortes et intrusive pour la sphère privée des étudiants, ces derniers se sentant souvent épiés par un observateur humain inconnu. Pour pallier ces contraintes de scalabilité, une seconde génération de systèmes de télésurveillance automatisés a vu le jour, s'appuyant sur des algorithmes heuristiques pour détecter des anomalies de base, telles que l'absence d'un visage ou la détection d'un double écran. Néanmoins, ces systèmes heuristiques souffrent d'un taux de faux positifs excessivement élevé, générant un stress cognitif indu chez l'apprenant qui craint d'être injustement accusé de fraude pour de simples mouvements oculaires naturels ou des bruits ambiants inoffensifs.

C'est dans ce contexte de double crise, marquée par la rigidité pédagogique et les limites technologiques de la télésurveillance, que survient la révolution contemporaine de l'intelligence artificielle, et plus spécifiquement l'avènement spectaculaire des grands modèles de langage. Ces réseaux de neurones profonds, reposant sur l'architecture Transformer introduite en 2017, ont démontré des capacités émergentes de compréhension, de raisonnement et de génération de texte naturel d'une fluidité et d'une pertinence sans précédent. Entraînés sur des corpus textuels gigantesques englobant une vaste proportion de la connaissance humaine numérisée, des modèles tels que LLaMA ou Mistral possèdent une compréhension encyclopédique de multiples domaines scientifiques. Leur intégration dans le secteur de l'éducation technologique marque la transition d'un paradigme de gestion de l'apprentissage vers un paradigme d'apprentissage agentique. L'intelligence artificielle agentique se distingue par sa capacité non seulement à traiter de l'information de manière passive, mais à percevoir son environnement, à poursuivre des objectifs complexes, à raisonner sur des plans d'action et à exécuter des tâches de manière autonome en interagissant avec d'autres systèmes. Dans le domaine pédagogique, un agent intelligent peut analyser l'historique d'apprentissage d'un étudiant, repérer avec précision ses lacunes conceptuelles et générer dynamiquement un parcours de remédiation sur mesure. Il peut agir comme un tuteur socratique disponible en permanence, capable de répondre aux interrogations de l'étudiant en s'appuyant exclusivement sur les ressources validées du cours grâce à des techniques d'ancrage sémantique. De surcroît, l'intelligence artificielle générative ouvre la voie à l'automatisation de tâches d'ingénierie pédagogique fastidieuses, telles que la conception de questionnaires complexes et la correction détaillée d'évaluations ouvertes, une tâche jusqu'alors exclusivement réservée aux enseignants en raison de la nuance sémantique requise.

Parallèlement, les avancées fulgurantes dans le domaine de la vision par ordinateur, portées par des architectures convolutives optimisées telles que YOLO, permettent d'envisager une nouvelle ère pour la télésurveillance des examens. Contrairement aux approches heuristiques rigides, les modèles de détection d'objets modernes peuvent analyser les flux vidéo en temps réel pour identifier avec une haute précision la présence d'objets contondants tels que des téléphones portables ou des écouteurs, tout en analysant l'orientation du visage et du regard avec une grande robustesse aux variations d'éclairage. L'intégration de ces technologies visuelles avec des modules d'analyse comportementale qui scrutent les interactions de l'étudiant avec l'interface numérique (frappes au clavier, mouvements de souris, changements de focus) permet de construire un score d'intégrité holistique. Ce score, affiné par un agent d'intelligence artificielle dédié, ne se contente plus de lever des alertes binaires, mais produit un rapport argumenté distinguant les comportements véritablement frauduleux des anomalies bénignes, allégeant ainsi considérablement la charge d'examen humain post-session. C'est à l'intersection de toutes ces innovations de rupture que se positionne le projet Fraudly. Né de la volonté d'offrir une solution complète et intégrée aux défis contemporains de l'enseignement supérieur, particulièrement ressentis au sein des institutions d'ingénierie telles que l'ENSET Mohammedia, Fraudly ambitionne de redéfinir l'écosystème numérique d'apprentissage. En combinant un moteur d'apprentissage adaptatif soutenu par des modèles de langage de pointe et un module de télésurveillance intelligent exploitant la vision par ordinateur, la plateforme se propose de réconcilier la personnalisation extrême de la pédagogie avec la garantie inébranlable de l'intégrité académique.
